% !TEX program = xelatex
\documentclass[11pt,a4paper]{article}
\usepackage{algo-preamble}

\title{\textbf{L16 \textperiodcentered{} Δυναμικός Προγραμματισμός III}\\[2pt]
\large Ευθυγράμμιση Συμβολοσειρών}
\author{Algorithms Class Hub}
\date{\small Αλγόριθμοι και Πολυπλοκότητα (K17) \textperiodcentered{} ΕΚΠΑ}

\begin{document}
\maketitle

\begin{center}\itshape\small
Ομοιότητα συμβολοσειρών και απόσταση επεξεργασίας (edit distance), ο αλγόριθμος
ευθυγράμμισης ως συντομότερο μονοπάτι σε πλέγμα, και η ευθυγράμμιση σε γραμμικό
χώρο κατά Hirschberg.
\end{center}

\section{Τι θα δούμε}

Πόσο \textbf{όμοιες} είναι δύο συμβολοσειρές; Το ερώτημα ακούγεται απλό, αλλά
είναι η βάση για το \texttt{unix diff}, για την αναγνώριση φωνής, και για τη
σύγκριση ακολουθιών DNA. Η απάντηση --- η \textbf{απόσταση επεξεργασίας} ---
είναι ένα κομψό DP. Και θα δούμε δύο βαθιές ιδέες:
\begin{enumerate}
  \item Το DP πλέγμα είναι στην πραγματικότητα \textbf{γράφημα} --- και η
  βέλτιστη λύση είναι ένα \textbf{συντομότερο μονοπάτι}.
  \item Με συνδυασμό DP \textbf{και} διαίρει-και-κυρίευε, μπορούμε να λύσουμε το
  πρόβλημα σε \textbf{γραμμικό χώρο} (Hirschberg).
\end{enumerate}

\section{Ομοιότητα συμβολοσειρών}

Δες τις δύο λέξεις:

\begin{codebox}
\begin{verbatim}
o c u r r a n c e
o c c u r e n c e
\end{verbatim}
\end{codebox}

Είναι ολοφάνερα «η ίδια λέξη με λάθη». Αλλά πώς το \textbf{ποσοτικοποιείς}; Η
ιδέα: μέτρα πόσες \textbf{στοιχειώδεις αλλαγές} χρειάζονται για να μετατρέψεις τη
μία στην άλλη.

\begin{keypoint}
\textbf{Απόσταση επεξεργασίας} (edit distance) --- [Levenshtein 1966,
Needleman-Wunsch 1970]. Δύο είδη «κόστους»:
\begin{itemize}
  \item \textbf{κόστος κενού} $\delta$ --- αφήνεις έναν χαρακτήρα
  \textbf{αταίριαστο} (μια εισαγωγή ή διαγραφή)·
  \item \textbf{κόστος σύγκρουσης} $\alpha_{xy}$ --- ταιριάζεις δύο
  \textbf{διαφορετικούς} χαρακτήρες $x \ne y$ (μια αντικατάσταση). Αν $x = y$,
  το κόστος είναι $0$.
\end{itemize}
Συνολικό κόστος $=$ άθροισμα όλων των κενών και των συγκρούσεων.
\end{keypoint}

\section{Ευθυγράμμιση συμβολοσειρών}

\begin{keypoint}
\textbf{Ευθυγράμμιση} (alignment) $M$ δύο συμβολοσειρών $X = x_1 \cdots x_m$ και
$Y = y_1 \cdots y_n$: ένα σύνολο διατεταγμένων ζευγών $x_i - y_j$ τέτοιο ώστε:
\begin{itemize}
  \item κάθε χαρακτήρας ανήκει σε \textbf{το πολύ ένα} ζεύγος, και
  \item τα ζεύγη \textbf{δεν διασταυρώνονται} --- τα $x_i - y_j$ και
  $x_{i'} - y_{j'}$ διασταυρώνονται αν $i < i'$ αλλά $j > j'$.
\end{itemize}
\[
\text{κόστος}(M) = \sum_{(x_i, y_j) \in M,\ x_i \ne y_j} \alpha_{x_i y_j} \ +
\sum_{x_i \text{ αταίριαστο}} \delta \ + \sum_{y_j \text{ αταίριαστο}} \delta
\]
\end{keypoint}

\begin{intuition}
\textbf{«Δεν διασταυρώνονται»} σημαίνει ότι αν διαβάσεις και τις δύο λέξεις από
αριστερά προς δεξιά, τα ζεύγη διατηρούν τη σειρά τους --- δεν μπερδεύεσαι.
\textbf{Στόχος:} βρες την ευθυγράμμιση \textbf{ελάχιστου κόστους}.
\end{intuition}

\section{Η αναδρομή}

Ορίζουμε $\text{OPT}(i, j) = $ ελάχιστο κόστος ευθυγράμμισης των προθεμάτων
$x_1 \cdots x_i$ και $y_1 \cdots y_j$. Κοιτάμε τους τελευταίους χαρακτήρες
$x_i, y_j$ --- υπάρχουν \textbf{τρεις} δυνατότητες:
\begin{itemize}
  \item \textbf{Ταίριασμα} $x_i$ με $y_j$: πληρώνεις $\alpha_{x_i y_j}$ (ή $0$
  αν ίδιοι) $+$ το βέλτιστο για $x_1 \cdots x_{i-1}$, $y_1 \cdots y_{j-1}$.
  \item \textbf{Το $x_i$ αταίριαστο:} πληρώνεις κενό $\delta$ $+$ το βέλτιστο
  για $x_1 \cdots x_{i-1}$, $y_1 \cdots y_j$.
  \item \textbf{Το $y_j$ αταίριαστο:} πληρώνεις κενό $\delta$ $+$ το βέλτιστο
  για $x_1 \cdots x_i$, $y_1 \cdots y_{j-1}$.
\end{itemize}

\begin{keypoint}
\[
\text{OPT}(i, j) =
\begin{cases}
j \cdot \delta & i = 0 \\[3pt]
i \cdot \delta & j = 0 \\[3pt]
\min \begin{cases}
\alpha_{x_i y_j} + \text{OPT}(i-1, j-1) \\
\delta + \text{OPT}(i-1, j) \\
\delta + \text{OPT}(i, j-1)
\end{cases} & \text{αλλιώς}
\end{cases}
\]
Οι βασικές περιπτώσεις: ευθυγράμμιση με την \textbf{κενή} συμβολοσειρά κοστίζει
ένα κενό ανά χαρακτήρα.
\end{keypoint}

\begin{codebox}
\begin{verbatim}
Ευθυγράμμιση(m, n, x, y, δ, α):
  Για i = 0 έως m:  M[i, 0] <- i*δ
  Για j = 0 έως n:  M[0, j] <- j*δ
  Για i = 1 έως m:
    Για j = 1 έως n:
      M[i, j] <- min( α(x_i, y_j) + M[i-1, j-1],
                      δ + M[i-1, j],
                      δ + M[i, j-1] )
  επίστρεψε M[m, n]
\end{verbatim}
\end{codebox}

\textbf{Παράδειγμα.} $x = \texttt{GAC}$, $y = \texttt{AGC}$, με $\delta = 1$ και
$\alpha_{xy} = 1$ για $x \ne y$:

\begin{center}
\begin{tabular}{@{}l *{4}{c}@{}}
\toprule
 & $\emptyset$ & A & G & C \\
\midrule
\textbf{$\emptyset$} & 0 & 1 & 2 & 3 \\
\textbf{G}           & 1 & 1 & 1 & 2 \\
\textbf{A}           & 2 & 1 & 2 & 2 \\
\textbf{C}           & 3 & 2 & 2 & \textbf{2} \\
\bottomrule
\end{tabular}
\end{center}

Η απόσταση επεξεργασίας είναι $M[3, 3] = \mathbf{2}$.

\textbf{Πολυπλοκότητα:} $\Theta(mn)$ χρόνος \textbf{και} $\Theta(mn)$ χώρος.

\section{Το DP πλέγμα είναι γράφημα}

Εδώ έρχεται η βαθιά ιδέα. Φτιάχνουμε ένα \textbf{γράφημα απόστασης επεξεργασίας}:
μια κορυφή για κάθε ζεύγος $(i, j)$, με ακμές:
\begin{itemize}
  \item \textbf{δεξιά} $(i, j-1) \to (i, j)$ --- κόστος $\delta$ (κενό στο $X$),
  \item \textbf{κάτω} $(i-1, j) \to (i, j)$ --- κόστος $\delta$ (κενό στο $Y$),
  \item \textbf{διαγώνια} $(i-1, j-1) \to (i, j)$ --- κόστος $\alpha_{x_i y_j}$
  (ταίριασμα).
\end{itemize}

\begin{keypoint}
Με αυτές τις ακμές, το $\text{OPT}(i, j)$ ισούται \textbf{ακριβώς} με το μήκος
του \textbf{συντομότερου μονοπατιού} από την κορυφή $(0, 0)$ στην $(i, j)$. Η
απόσταση επεξεργασίας είναι το συντομότερο μονοπάτι $(0,0) \to (m, n)$.
\end{keypoint}

\section{Ευθυγράμμιση σε γραμμικό χώρο --- Hirschberg}

Ο πίνακας $\Theta(mn)$ είναι το πρόβλημα: στην υπολογιστική βιολογία
$m = n = 100.000$, και ένας πίνακας $10^{10}$ κελιών \textbf{δεν χωράει} στη
μνήμη.

\subsection{Πρώτη παρατήρηση --- η τιμή σε γραμμικό χώρο}

Για να υπολογίσεις τη γραμμή $\text{OPT}(i, \cdot)$ χρειάζεσαι \textbf{μόνο} την
προηγούμενη γραμμή $\text{OPT}(i-1, \cdot)$. Άρα η \textbf{τιμή}
$\text{OPT}(m, n)$ υπολογίζεται σε $O(m + n)$ χώρο και $O(mn)$ χρόνο, κρατώντας
δύο γραμμές κάθε φορά.

\begin{warningbox}
\textbf{Αλλά} --- έτσι χάνεις τον υπόλοιπο πίνακα, και η \textbf{ίδια η
ευθυγράμμιση} δεν ανακτάται με το συνηθισμένο πέρασμα προς τα πίσω. Χρειάζεται
έξυπνη ιδέα.
\end{warningbox}

\subsection{Η ιδέα του Hirschberg (1975) --- διαίρει και κυρίευε $+$ DP}

Δουλεύουμε στο γράφημα απόστασης. Ορίζουμε:
\begin{itemize}
  \item $f(i, j) = $ συντομότερο μονοπάτι από $(0, 0)$ στο $(i, j)$ --- δηλαδή
  $\text{OPT}(i, j)$·
  \item $g(i, j) = $ συντομότερο μονοπάτι από $(i, j)$ στο $(m, n)$ ---
  υπολογίζεται \textbf{όμοια}, αντιστρέφοντας τις ακμές.
\end{itemize}

\begin{keypoint}
Το κόστος του συντομότερου μονοπατιού $(0,0) \to (m,n)$ \textbf{που περνά} από
την κορυφή $(i, j)$ είναι ακριβώς $f(i, j) + g(i, j)$.
\end{keypoint}

Διαλέγουμε τη \textbf{μεσαία στήλη} $j = n/2$. Έστω $q$ ο δείκτης γραμμής που
ελαχιστοποιεί την ποσότητα $f(q,\, n/2) + g(q,\, n/2)$. Τότε \textbf{υπάρχει}
βέλτιστη ευθυγράμμιση που περνά από την κορυφή $(q,\, n/2)$.

\textbf{Διαίρει:} βρες το $q$ (με DP σε γραμμικό χώρο, υπολογίζοντας τη στήλη
$f(\cdot, n/2)$ και $g(\cdot, n/2)$). Ευθυγράμμισε το $(q, n/2)$.

\textbf{Κυρίευε:} λύσε αναδρομικά τα \textbf{δύο μισά} --- το πάνω-αριστερά
ορθογώνιο $(0,0)$--$(q, n/2)$ και το κάτω-δεξιά $(q, n/2)$--$(m, n)$.

\begin{keypoint}
\textbf{Πολυπλοκότητα:} $O(m + n)$ \textbf{χώρος} και $O(mn)$ \textbf{χρόνος}. Ο
χρόνος μένει $O(mn)$ γιατί κάθε επίπεδο της αναδρομής δουλεύει σε ορθογώνια
συνολικού εμβαδού που \textbf{υποδιπλασιάζεται} ---
$mn + mn/2 + mn/4 + \dots = O(mn)$.
\end{keypoint}

\begin{intuition}
\textbf{Γιατί είναι όμορφο:} συνδυάζει \textbf{δύο} από τα μοντέλα του
μαθήματος. Το DP βρίσκει τη μεσαία κορυφή· το διαίρει-και-κυρίευε σπάει το
πρόβλημα γύρω της. Κανένα από τα δύο μόνο του δεν θα έδινε γραμμικό χώρο.
\end{intuition}

\begin{recapbox}
\textbf{Τι κρατάμε από αυτή τη διάλεξη}
\begin{enumerate}
  \item \textbf{Απόσταση επεξεργασίας} --- ελάχιστο κόστος μετατροπής μιας
  συμβολοσειράς σε άλλη, με κόστος κενού $\delta$ και κόστος σύγκρουσης $\alpha$.
  \item \textbf{Ευθυγράμμιση} $=$ μη-διασταυρούμενα ζεύγη χαρακτήρων. Αναδρομή με
  \textbf{τρεις} περιπτώσεις: ταίριασμα (διαγώνια), κενό στο $X$ (κάτω), κενό στο
  $Y$ (δεξιά). Χρόνος \& χώρος $\Theta(mn)$.
  \item Το DP πλέγμα είναι \textbf{γράφημα}: $\text{OPT}(i,j) = $ συντομότερο
  μονοπάτι $(0,0) \to (i,j)$.
  \item \textbf{Hirschberg} --- συνδυασμός DP $+$ διαίρει-και-κυρίευε. Με $f$
  (εμπρός) και $g$ (πίσω), βρες την κορυφή της μεσαίας στήλης που ελαχιστοποιεί
  $f + g$, και αναδρομή. $O(m+n)$ χώρος, $O(mn)$ χρόνος.
\end{enumerate}
\end{recapbox}

\lecturefooter
\end{document}
